\section{Discussion and Conclusion}
% =====================================================================

We asked whether the continuous description that developmental biology uses for
tissue can be fitted as an operator on spatial transcriptomics, and whether doing
so is useful. On the evidence presented, it is. Across \MSSections{} sections and
four technologies, treating tissue as a field rather than a graph improves
prediction at unmeasured coordinates against every baseline tested, with
Holm-corrected significance and large effect sizes under paired tests. A
parameter-matched control attributes the gain to the operator rather than to
capacity; a non-spatial control and an implicit neural field together show that
neither message passing nor continuity alone accounts for it; the reconstruction
preserves the tissue-level structure a downstream analysis depends on; and replacing the
exponential integrator costs $\NumSpeedup\times$ more steps to reach the same
tolerance, so the numerical machinery is load-bearing.

We also delimit what this class of method can establish:
Theorem~\ref{thm:identifiability} shows a stationary snapshot constrains the
reaction network only on a null set, which both explains our perturbation result
and identifies the data needed to go further. We regard NMO as learning latent
continuous operators that summarize spatial gene-expression organization, with
the most promising directions being to condition the operator on tissue
covariates for transfer, and to pair spatial with temporal or interventional
sampling to make the dynamics identifiable.

\label{endofmain}
